\begin{abstract}
Aerial sensing can reveal regions that a ground robot cannot readily inspect,
but maximizing environmental information alone need not reveal the locations
that make a manipulation task executable. We study this mismatch in cooperative
retrieval by an aerial robot and a ground mobile manipulator. Our framework
projects validated inverse-reachability candidates through the ground robot's
footprint to identify environment cells relevant to future manipulation. It
combines this support with an observation-deficit field and a finite-window
LiDAR acquisition model to select aerial viewpoints. A shared execution interface
requires measured ground support, geometric compatibility and collision-aware
manipulation screening before handoff; sensing predictions never count as
observations. We evaluate the complete perception-to-retrieval loop on a
P450--BUNKER--AUBO i5--AG95 physics-based simulation platform using 96 independent
scenes and 192 primary paired trials. Under the same three-window observation
budget, manipulation-aware selection achieves physical retrieval in 80/96
scenes, compared with 56/96 for a generic NBV policy that differs only in gain
weighting. The paired difference is 25.00 percentage points, with a prespecified
conservative 95\% confidence interval of [8.56, 39.15] and an exact McNemar
$p=8.43\times10^{-6}$. These results support task-conditioned sensing for the
studied static, known-object retrieval population; they do not establish a
general time-saving or real-robot performance claim.
\end{abstract}
